\documentclass[11pt]{article}

\usepackage[a4paper,margin=0.82in]{geometry}
\usepackage[T1]{fontenc}
\usepackage[utf8]{inputenc}
\usepackage{lmodern}
\usepackage{amsmath,amssymb,amsthm,mathtools}
\usepackage{booktabs,enumitem,microtype,tabularx,xcolor}
\usepackage[numbers,sort&compress]{natbib}
\usepackage[colorlinks=true,linkcolor=blue!65!black,
 citecolor=blue!65!black,urlcolor=blue!70!black]{hyperref}
\usepackage[nameinlink,noabbrev]{cleveref}
\hypersetup{
 pdftitle={The TPC-18 Full-Block Physical Remainder},
 pdfauthor={Liang Wang},
 pdfsubject={TPC-117 exact cover, canonical residual, and dual stop certificate},
 pdfkeywords={full-block cover, physical remainder, Moore-Penrose inverse, dual certificate, fixed shift}
}

\newtheorem{theorem}{Theorem}[section]
\newtheorem{proposition}[theorem]{Proposition}
\newtheorem{corollary}[theorem]{Corollary}
\newtheorem{lemma}[theorem]{Lemma}
\theoremstyle{definition}
\newtheorem{definition}[theorem]{Definition}
\theoremstyle{remark}
\newtheorem{remark}[theorem]{Remark}

\newcommand{\C}{\mathbb C}
\newcommand{\one}{\mathbf 1}
\newcommand{\Lzero}{\mathrm{L0}}
\newcommand{\Lone}{\mathrm{L1}}
\newcommand{\Ltwo}{\mathrm{L2}}
\newcommand{\Ran}{\operatorname{ran}}
\newcommand{\Ker}{\operatorname{ker}}
\newcommand{\Block}{\mathcal B}

\title{\textbf{The TPC-18 Full-Block Physical Remainder:}\\
Exact Range Tests, Canonical Residuals,\\
and Literal Dual Stop Certificates}
\author{Liang Wang\\
School of Artificial Intelligence and Automation\\
Huazhong University of Science and Technology, Wuhan 430074, P.R. China\\
\texttt{wangliang.f@gmail.com}}
\date{July 2026}

\begin{document}
\maketitle

\begin{abstract}
TPC-18 derives a powerful determinant-dispersion estimate on
admissible hyperbola blocks, but its global fixed-shift conclusion is
conditional on those blocks covering the full physical tail, or on
the uncovered physical remainder being \(o(X)\).  We isolate that
condition as a finite exact range problem.

Let \(w_X\) be the literal physical coefficient vector on the
complete active atom set, and let
\[
 B_X:\C^{\Block_X}\longrightarrow\C^{\Omega_X}
\]
be the synthesis whose columns are the actual TPC-18 blocks with
their inherited overlap weights, signs, cutoffs and normalization.
Writing \(P_X=B_XB_X^\dagger\), we prove
\[
 w_X\in\Ran B_X
 \iff (I-P_X)w_X=0
 \iff \langle y,w_X\rangle=0
       \quad\text{for all }y\in\Ker B_X^*.
\]
The canonical coefficient and residual are
\[
 c_X^*=B_X^\dagger w_X,\qquad
 r_X^*=(I-P_X)w_X,
\]
and the physical scalar splits exactly into a block term plus
\(\nu_X\langle a_X,r_X^*\rangle\).  The sharp distance to the block
carrier is \(\|r_X^*\|\); when it is nonzero, a maximizing unit dual
witness lies in \(\Ker B_X^*\).  Thus one nonzero exact witness disproves exact
cover, while a nonzero residual is still compatible with the route
only if its literal physical pairing is separately \(o(X)\).

Exact cover alone is not an arithmetic estimate: transporting a
block theorem costs the actual canonical coefficients or the
relevant singular-value bound, and that cost enters the physical
ledger.  The present repository does not contain a machine-readable
growing TPC-18 block matrix and coefficient vector from which the
range test could be executed.  Accordingly the full-block
hypothesis remains open.  The result is an \(\Lzero/\Lone\) cover
audit, not an \(\Ltwo\) fixed-shift saving or a twin-prime theorem.
\end{abstract}

\noindent\textbf{Keywords:}
full-block reassembly; exact cover; physical residual; dual witness;
least squares; endpoint ledger.

\medskip
\noindent\textbf{Literal-carrier convention.}
The row set is required to be the complete physical atom set
specified by a verified TPC-116 archive.  A
column is an actual admissible TPC-18 block, not an abstract block of
similar dimensions.  Overlap weights, both polarizations, cutoffs,
prime-power errors, the fixed \(h_0\ne0\), and the global
normalization are part of the matrix data.

\section{The inherited gap}

TPC-18 constructs a Type-II tail interface and proves conditional
block estimates under a determinant-dispersion hypothesis
\citep{WangTPC18}.  Its conditional fixed-shift conclusion assumes
an additional global all-block identity
\[
 H_{h_0,R;V}^{\mathrm{HB}}(X)
 =\sum_{b\in\Block_X}\bigl(C_b+T_b\bigr)+o(X),
\]
including every geometric and normalization requirement.  The paper
does not prove that global cover.  TPC-92 later stresses that an
exact full-block cover is sufficient, while any nonzero residual
requires a separate physical theorem \citep{WangTPC92}.  TPC-116
returns the same issue after completing the complement taxonomy
\citep{WangTPC116}.

The word ``cover'' can mean three different things:
\begin{enumerate}[leftmargin=2em]
\item every atom belongs to at least one set-theoretic block;
\item the incidence columns span the coefficient support;
\item the \emph{weighted signed physical vector} belongs to the
  range of the literal block synthesis.
\end{enumerate}
Only the third statement supplies the identity needed by TPC-18.
The first two can fail after overlap weights, masks and signs are
restored.

\section{The literal finite model}

Let \(\Omega\) be the finite active atom set and equip
\(\mathcal H=\C^\Omega\) with its declared physical inner product.
Let \(w\in\mathcal H\) be the literal coefficient vector.  For the
finite block set \(\Block\), define
\[
 B:\C^\Block\to\mathcal H,\qquad
 Bc=\sum_{b\in\Block}c_b\,v_b,
\]
where \(v_b\) contains every atom weight contributed by block \(b\).
The matrix is allowed to have repeated, dependent or zero columns.
Every adjoint, orthogonal projection, norm and Moore--Penrose inverse
below is computed relative to this declared physical inner product
and the declared coefficient-space inner product; a hidden change to
an unweighted Euclidean geometry changes the certificate.
Put
\[
 P=BB^\dagger,\qquad
 c^*=B^\dagger w,\qquad
 r^*=(I-P)w.
\]

\begin{theorem}[Exact cover alternatives]
\label{thm:cover}
The following are equivalent:
\begin{enumerate}[label=(\roman*),leftmargin=2em]
\item \(w\in\Ran B\);
\item \(Pw=w\);
\item \(r^*=0\);
\item \(\langle y,w\rangle=0\) for every \(y\in\Ker B^*\).
\end{enumerate}
When they hold, \(c^*\) is the unique least-norm exact cover:
\[
 Bc^*=w,\qquad
 \|c^*\|_2=\min_{Bc=w}\|c\|_2.
\]
\end{theorem}

\begin{proof}
The Moore--Penrose identity makes \(P=BB^\dagger\) the orthogonal
projection onto \(\Ran B\).  Hence (i)--(iii) are equivalent.
Finite-dimensional orthogonality gives
\((\Ran B)^\perp=\Ker B^*\), proving (iv).  The least-norm property of
\(B^\dagger w\) follows by decomposing any solution into
\((\Ker B)^\perp\oplus\Ker B\).
\end{proof}

\begin{proposition}[Sharp residual and dual certificate]
\label{prop:dual}
For every \(w\),
\[
 \boxed{\quad
 \operatorname{dist}(w,\Ran B)=\|r^*\|_2
 =\max_{\substack{y\in\Ker B^*\\\|y\|_2\le1}}
 |\langle y,w\rangle|.
 \quad}
\]
If \(r^*\ne0\), the maximum is attained by
\(y=r^*/\|r^*\|_2\).  Thus any exact vector
\(y\in\Ker B^*\) with \(\langle y,w\rangle\ne0\) is a literal
certificate that the proposed block family does not cover \(w\).
\end{proposition}

\begin{proof}
Orthogonal projection gives the distance identity.  For
\(y\in\Ker B^*=(\Ran B)^\perp\),
\(\langle y,w\rangle=\langle y,r^*\rangle\), so
Cauchy--Schwarz gives the upper bound and the stated vector attains
it when \(r^*\ne0\).  When \(r^*=0\), the vector \(y=0\) belongs to
the closed unit ball and attains the value zero.
\end{proof}

\begin{remark}[A range witness is stronger than a missing-set label]
A set-theoretically uncovered atom gives a coordinate witness only
when no block column has any component in that coordinate.  In
general, signed overlap can create subtler cokernel vectors.
\Cref{prop:dual} detects both mechanisms without changing the
physical weights.
\end{remark}

\section{The physical scalar and the allowed remainder}

Let \(a\in\mathcal H\) be the physical evaluation vector and let
\(\nu_X\) be the inherited scalar normalization.  The target is
\[
 H_X=\nu_X\langle a,w\rangle.
\]

\begin{theorem}[Exact block--remainder identity]
\label{thm:physical}
For the canonical pair \(c^*,r^*\),
\[
 \boxed{\quad
 H_X
 =\nu_X\langle B^*a,c^*\rangle
  +\nu_X\langle a,r^*\rangle.
 \quad}
\]
Consequently the TPC-18 cover gate closes if either:
\begin{enumerate}[label=(\alph*),leftmargin=2em]
\item \(r^*=0\) and the complete block term is \(o(X)\); or
\item the complete block term is \(o(X)\) and the literal physical
  residual satisfies
  \(\nu_X\langle a,r^*\rangle=o(X)\).
\end{enumerate}
No smallness of \(\|r^*\|_2\) is necessary if the physical pairing is
proved small; conversely, \(r^*\ne0\) alone does not prove that its
physical pairing is large.
\end{theorem}

\begin{proof}
Since \(w=Bc^*+r^*\), take its inner product with \(a\) and use
adjointness.  The two sufficient alternatives are immediate.
\end{proof}

\begin{corollary}[A norm-based sufficient residual estimate]
If
\[
 \nu_X\|a\|_2\|r^*\|_2=o(X),
\]
then the physical residual in \cref{thm:physical} is \(o(X)\).
This criterion is sufficient and sharp given only the two norms.
\end{corollary}

\begin{proof}
Cauchy--Schwarz proves the bound; collinear vectors show sharpness.
\end{proof}

\section{Exact cover is not an estimate}

Suppose a block theorem bounds the vector of block evaluations
\[
 q=B^*a,\qquad |q_b|\le U_b.
\]
Even under exact cover,
\[
 |H_X|
 \le\nu_X\sum_{b\in\Block}|c_b^*|U_b.
\]
Thus the number, size and conditioning of the canonical coefficients
matter.

\begin{proposition}[Coefficient transport costs]
\label{prop:transport}
If \(r^*=0\), then
\[
 |H_X|
 \le\nu_X\|c^*\|_2\|B^*a\|_2
 \le\nu_X\|B^\dagger\|\,\|w\|_2\,\|B^*a\|_2.
\]
Both inequalities have optimal constant one.  If instead individual
block bounds are inserted, the exact triangle cost is
\(\sum_b|c_b^*|U_b\).  Neither cost may be omitted merely because
the cover is exact.
\end{proposition}

\begin{proof}
Use Cauchy--Schwarz, then
\(\|c^*\|=\|B^\dagger w\|\le\|B^\dagger\|\|w\|\).
Equality examples are supplied by singular vectors and collinear
evaluation vectors.
\end{proof}

\begin{remark}[Relation to TPC-113]
TPC-113 identifies the quotient condition number of a literal
physical synthesis \citep{WangTPC113}.  Here \(B\) is a different,
specific full-block synthesis.  A previously paid frame cost may be
reused only through an exact factorization identifying the same
operator event; otherwise \(\|B^\dagger\|\) is a new ledger input.
\end{remark}

\section{What an executable TPC-18 archive must contain}

An actual range decision requires:
\begin{enumerate}[leftmargin=2em]
\item a stable row identifier for every physical atom returned by
  the complete TPC-116 partition;
\item a stable column identifier for every \(C_b\) and \(T_b\), with
  the allowed \(L,K,I\) and primitive/nonprimitive status;
\item exact or certified interval entries including overlap weights,
  cutoffs, both polarizations, M\"obius/von Mangoldt factors and the
  prime-power correction;
\item the literal vector \(w_X\), evaluation \(a_X\), global
  normalization \(\nu_X\), and fixed \(h_0\);
\item a certificate \(Bc=w\), or a dual witness
  \(B^*y=0\), \(\langle y,w\rangle\ne0\), or the canonical residual
  plus a quantitative physical estimate; and
\item the coefficient/conditioning cost used to transport the block
  dispersion theorem.
\end{enumerate}
Dimensions and set-theoretic coverage statistics are useful
diagnostics but do not replace these fields.

\section{Finite exact regression certificate}

The script
\texttt{experiments/tpc117\_cover\_certificate.py} uses
\texttt{fractions.Fraction} to verify on an overlapping block matrix:
\begin{enumerate}[leftmargin=2em]
\item an exact cover \(Bc=w\);
\item an exact cokernel vector \(B^*y=0\);
\item a perturbed vector with nonzero dual pairing;
\item the canonical rational residual and distance identity,
  including the full-row-rank case \(\Ker B^*=\{0\}\); and
\item the exact physical block--remainder decomposition.
\end{enumerate}
The example demonstrates the certificate format only.  It is not a
sample of the growing TPC-18 carrier.

\section{Stop/go verdict and claim boundary}

\begin{center}
\small
\begin{tabularx}{\textwidth}{@{}p{0.19\textwidth}p{0.38\textwidth}X@{}}
\toprule
Status & Certificate & Consequence\\
\midrule
Go exact &
Literal archive and verified \(B_Xc_X=w_X\), with affordable
transport cost &
No physical cover residual; pass the charged block estimate to the
endpoint ledger.\\
Go with remainder &
Canonical identity plus
\(\nu_X\langle a_X,r_X^*\rangle=o(X)\) &
The nonzero residual is separately controlled at the original
normalization.\\
Stop the exact-cover subroute &
Exact \(B_X^*y_X=0\) and
\(\langle y_X,w_X\rangle\ne0\), with no residual theorem &
Exact cover is false; enlarge/change the carrier or prove the
physical residual small.\\
Current TPC state &
TPC-18 states the global cover as an additional hypothesis; no
machine-readable growing matrix/vector certificate is available
here &
H6 is mathematically OPEN; its present range test is
NOT-TESTABLE-FROM-CURRENT-ARTIFACTS, and its ledger cost is unknown.\\
\bottomrule
\end{tabularx}
\end{center}

\paragraph{Established.}
\Cref{thm:cover,prop:dual,thm:physical,prop:transport} are exact
\(\Lzero\) finite-dimensional theorems.  The required literal
TPC-18 archive schema and original-scale pairing are \(\Lone\)
interfaces.

\paragraph{Not established.}
We do not prove that the actual TPC-18 blocks cover the growing
physical vector, and we do not bound their canonical residual or
transport cost.  No determinant-dispersion hypothesis is upgraded.
There is no new \(\Ltwo\) fixed-\(h_0\) estimate, parity
breakthrough, Hardy--Littlewood asymptotic, prime-pair lower bound or
twin-prime theorem.  The next gate is the nonduplicated strict
\(1/400\) endpoint ledger.

\begingroup
\small
\setlength{\bibsep}{2pt plus 0.3ex}
\sloppy
\bibliographystyle{plainnat}
\bibliography{references}
\endgroup
\end{document}
